% Lean: EconLean/Behavioral/{PresentBias,RankDependent,ReferenceDependence,Temptation,Logit,RationalInattention,Ambiguity,GilboaSchmeidler,StochasticChoice}.lean
\chapter{Behavioral Decision Theory}
\label{chap:behavioral}

Sources: Laibson (1997); O'Donoghue--Rabin (1999); Strotz (1956); Quiggin (1982); Kahneman--Tversky (1979);
Tversky--Kahneman (1992); MWG \S6.D; K\H{o}szegi--Rabin (2006, 2007); Gul--Pesendorfer (2001);
Anderson--de Palma--Thisse (1988); Sims (2003); Mat\v{e}jka--McKay (2015); Ellsberg (1961);
Gilboa--Schmeidler (1989). Review record: \texttt{Audit/REVIEW-behavioral.md}.

\section{Present bias}
\label{sec:beh-present-bias}

\begin{definition}[Quasi-hyperbolic discounting]
\label{def:beh-discount}
\lean{Econ.PresentBias.discount, Econ.PresentBias.presentValue, Econ.PresentBias.PreferenceReversal}
\leanok
Self $t$ weights a payoff $k$ periods ahead by $D(0)=1$, $D(k)=\beta\delta^k$ ($k\ge1$). A \emph{preference
reversal} between $a$ at date $s>0$ and $b$ at $s+1$: from date $0$, $b$ at $s+1$ is strictly preferred; at date
$s$, $a$ now is strictly preferred.
\end{definition}

\begin{theorem}[Exponential discounting is time consistent]
\label{thm:beh-exp-consistent}
\lean{Econ.PresentBias.exponential_timeConsistent, Econ.PresentBias.not_preferenceReversal_of_beta_one}
\leanok
\uses{def:beh-discount}
With $\beta=1$, $\delta>0$ and $t\le t'\le s,s'$, self $t$ and self $t'$ rank dated rewards identically; in
particular there are no reversals.
\end{theorem}
\begin{proof}
\leanok
$\delta^{s-t}=\delta^{t'-t}\delta^{s-t'}$; dividing both sides by $\delta^{t'-t}>0$ preserves strict inequality.
\end{proof}

\begin{theorem}[Reversals iff present bias]
\label{thm:beh-reversal-iff}
\lean{Econ.PresentBias.exists_preferenceReversal_iff}
\leanok
\uses{def:beh-discount}
For $\beta,\delta>0$, positive rewards exhibiting a reversal exist iff $\beta<1$.
\end{theorem}
\begin{proof}
\leanok
A reversal means $\beta\delta b<a<\delta b$. If $\beta\ge1$ and $b>0$ this is impossible. If $\beta<1$ take $b=1$,
$a=\delta(1+\beta)/2$, $s=1$.
\end{proof}

\begin{definition}[O'Donoghue--Rabin task]
\label{def:beh-task}
\lean{Econ.PresentBias.Task, Econ.PresentBias.Task.U, Econ.PresentBias.Task.NaiveDoes, Econ.PresentBias.Task.naiveDate, Econ.PresentBias.Task.tcDate, Econ.PresentBias.Task.sophFrom, Econ.PresentBias.Task.sophDate}
\leanok
A task is done once in periods $1,\dots,T$; doing it at $\tau$ yields reward $v_\tau$ and cost $c_\tau$, $\delta=1$.
Self $t$ values doing it now at $\beta v_t-c_t$ (immediate costs) or $v_t-\beta c_t$ (immediate rewards), and at
$\tau>t$ at $\beta(v_\tau-c_\tau)$. A naif does it at $t$ iff $t=T$ or now is weakly better than every later date;
$\tau_n$ is the first such $t$ and $\tau_{tc}$ is $\tau_n$ with $\beta=1$. A sophisticate's date is computed by
backward induction from $T$ (doing it when indifferent); $\tau_s$ is its value from period $1$.
\end{definition}

\begin{theorem}[Sophisticates finish no later than naifs]
\label{thm:beh-soph-le-naive}
\lean{Econ.PresentBias.Task.soph_le_naive}
\leanok
\uses{def:beh-task}
For every task and every $\beta$, $\tau_s\le\tau_n$.
\end{theorem}
\begin{proof}
\leanok
At $t=\tau_n$ the naif prefers now to every later date, in particular to the sophisticate's continuation date,
so the sophisticated plan from $t$ stops at $t$. By induction downward from $t$, the plan from period $1$ stops
no later than $t$.
\end{proof}

\begin{theorem}[Naifs procrastinate with immediate costs, preproperate with immediate rewards]
\label{thm:beh-naive-tc}
\lean{Econ.PresentBias.Task.tc_le_naive_of_immediateCosts, Econ.PresentBias.Task.naive_le_tc_of_immediateRewards}
\leanok
\uses{def:beh-task}
If costs are immediate, $c\ge0$ and $0<\beta\le1$, then $\tau_{tc}\le\tau_n$. If rewards are immediate, $v\ge0$ and
$0\le\beta\le1$, then $\tau_n\le\tau_{tc}$.
\end{theorem}
\begin{proof}
\leanok
$\tau_{tc}$ is the earliest period whose net value $w=v-c$ is at least every later one. Immediate costs: if the
naif acted at $t<\tau_{tc}$ there is $\tau>t$ with $w_\tau>w_t$, yet $\beta w_\tau\le\beta v_t-c_t\le\beta w_t$,
contradicting $\beta>0$. Immediate rewards: at $\tau_{tc}$, $\beta w_\tau\le\beta w_{\tau_{tc}}\le v-\beta c$, so
the naif acts there.
\end{proof}

\begin{theorem}[The paper's examples]
\label{thm:beh-or-examples}
\lean{Econ.PresentBias.Task.report_dates, Econ.PresentBias.Task.movie_dates}
\leanok
\uses{def:beh-task}
With $\beta=1/2$, $T=4$ and values $3,5,8,13$: immediate costs give $(\tau_{tc},\tau_s,\tau_n)=(1,2,4)$;
immediate rewards give $(4,1,3)$.
\end{theorem}
\begin{proof}
\leanok
Unfold the backward induction and the definition of the first naive date and compare numbers.
\end{proof}

\section{Rank-dependent utility and prospect theory}
\label{sec:beh-rdu}

\begin{definition}[Grid lotteries, FOSD, RDU, OPT, CPT]
\label{def:beh-lottery}
\lean{Econ.RankDependent.Lottery, Econ.RankDependent.G, Econ.RankDependent.F, Econ.RankDependent.eu, Econ.RankDependent.FOSD, Econ.RankDependent.rdu, Econ.RankDependent.opt, Econ.RankDependent.cpt}
\leanok
Outcomes are grid points $0<1<\dots<n$ (index order is the ranking). A lottery is a probability vector $p$ on the
grid; $G_p(k)=\Pr(\ge k)$, $F_p(k)=\Pr(\le k)$. $p$ FOSD $q$ iff $G_q\le G_p$. RDU is
$\sum_k[w(G_p(k))-w(G_p(k+1))]u_k$; OPT is $\sum_k\pi(p_k)v_k$; CPT with reference point $r$ is
$\sum_{k>r}w^+(G_p(k))(v_k-v_{k-1})-\sum_{k<r}w^-(F_p(k))(v_{k+1}-v_k)$.
\end{definition}

\begin{theorem}[FOSD iff higher EU for all increasing $u$; MWG 6.D.1]
\label{thm:beh-fosd-eu}
\lean{Econ.RankDependent.fosd_iff_eu, Econ.RankDependent.eu_eq_layer}
\leanok
\uses{def:beh-lottery}
\end{theorem}
\begin{proof}
\leanok
Abel summation gives $E_pu=u_0+\sum_{k<n}G_p(k+1)(u_{k+1}-u_k)$, which is monotone in $G_p$. Conversely the
indicator of $\{j\ge k\}$ has expectation $G_p(k)$.
\end{proof}

\begin{theorem}[RDU: Quiggin]
\label{thm:beh-rdu}
\lean{Econ.RankDependent.rdu_id, Econ.RankDependent.rdu_mono, Econ.RankDependent.rdu_strictMono, Econ.RankDependent.sum_decisionWeight}
\leanok
\uses{def:beh-lottery, thm:beh-fosd-eu}
With $w=\mathrm{id}$, RDU is EU. With $w$ increasing, $w(0)=0$, $w(1)=1$, decision weights are nonnegative and sum
to one, and RDU respects FOSD (strictly when $w,u$ are strictly increasing and dominance is strict somewhere).
\end{theorem}
\begin{proof}
\leanok
Abel summation puts RDU in layer form $u_0+\sum_{k<n}w(G_p(k+1))(u_{k+1}-u_k)$, monotone in each $G_p(k)$.
\end{proof}

\begin{theorem}[OPT violates FOSD]
\label{thm:beh-opt}
\lean{Econ.RankDependent.opt_violates_fosd}
\leanok
\uses{def:beh-lottery}
If $\pi(0)=0$ and $\pi(a+b)<\pi(a)+\pi(b)$ for some $a,b>0$, $a+b\le1$, there are a strictly increasing $v$ on a
four-point grid and lotteries $p$ FOSD $q$ with $\mathrm{OPT}(p)<\mathrm{OPT}(q)$.
\end{theorem}
\begin{proof}
\leanok
$p$ puts $a+b$ on the top outcome; $q$ splits it as $a$ and $b$ on two outcomes just below. With $v=(0,1,1+\varepsilon,1+2\varepsilon)$
the OPT difference is at least $D-\varepsilon M>0$ for $D=\pi(a)+\pi(b)-\pi(a+b)$ and $\varepsilon=D/(2M)$.
\end{proof}

\begin{theorem}[CPT]
\label{thm:beh-cpt}
\lean{Econ.RankDependent.cpt_mono, Econ.RankDependent.cpt_id, Econ.RankDependent.cpt_zero_eq_rdu, Econ.RankDependent.cpt_rejects_fair_coin, Econ.RankDependent.tk_rejects_fair_coin}
\leanok
\uses{def:beh-lottery, thm:beh-rdu}
CPT respects FOSD; without weighting and with $v_r=0$ it is $E_pv$; with no losses it is RDU; with loss aversion
$v(-x)<-v(x)$ and $0<w^+(1/2)\le w^-(1/2)$ it rejects the fair coin $\pm x$, in particular for the TK (1992)
value function with $\lambda>1$.
\end{theorem}
\begin{proof}
\leanok
Monotonicity term by term, using $F_p(k)=1-G_p(k+1)$. The EU identity splits the layer form at $r$ and
telescopes the loss part. The fair coin's CPT value is $w^+(\tfrac12)v_2+w^-(\tfrac12)v_0<0$.
\end{proof}

\section{Reference-dependent preferences}
\label{sec:beh-kr}

\begin{definition}[K\H{o}szegi--Rabin utility and equilibria]
\label{def:beh-kr}
\lean{Econ.ReferenceDependence.u, Econ.ReferenceDependence.U, Econ.ReferenceDependence.IsPE, Econ.ReferenceDependence.IsPPE, Econ.ReferenceDependence.IsCPE, Econ.ReferenceDependence.IsPEDet, Econ.ReferenceDependence.μLin}
\leanok
$u(c\mid r)=\sum_k m_k(c_k)+\sum_k\mu(m_k(c_k)-m_k(r_k))$; $U(F\mid G)$ averages it. $F$ is a PE if no $F'\in D$ does
better given expectation $F$; a PPE is a best PE; a CPE maximizes $U(F\mid F)$. Linear $\mu$: slope $\eta$ on gains,
$\eta\lambda$ on losses.
\end{definition}

\begin{theorem}[Existence and multiplicity]
\label{thm:beh-kr-exist}
\lean{Econ.ReferenceDependence.exists_ppe, Econ.ReferenceDependence.exists_cpe, Econ.ReferenceDependence.isPEDet_of_isMax, Econ.ReferenceDependence.exists_ppeDet, Econ.ReferenceDependence.two_PE_example}
\leanok
\uses{def:beh-kr}
On finite choice sets a PPE exists whenever a PE does and a CPE exists. In riskless choice with $\eta\ge0$,
$\lambda\ge1$ the consumption-utility maximizer is a PE. PE need not be unique.
\end{theorem}
\begin{proof}
\leanok
Maxima over finite sets. For the riskless PE, $\mu(x)\le\eta x$ gives $u(c'\mid c)\le(1+\eta)m(c')-\eta m(c)\le m(c)$.
\end{proof}

\begin{theorem}[Endowment effect]
\label{thm:beh-endowment}
\lean{Econ.ReferenceDependence.sell_iff, Econ.ReferenceDependence.buy_iff, Econ.ReferenceDependence.endowment_effect, Econ.ReferenceDependence.keep_isPE_iff, Econ.ReferenceDependence.sell_isPE_iff}
\leanok
\uses{def:beh-kr}
With the status quo as reference, $\mathrm{WTA}=v(1+\eta\lambda)/(1+\eta)$, $\mathrm{WTP}=v(1+\eta)/(1+\eta\lambda)$ and
$\mathrm{WTP}<\mathrm{WTA}$ iff $\lambda>1$. With expectations as reference, keeping is a PE iff $p\le\mathrm{WTA}$
and selling iff $p\ge v(1+\eta)/(1+\eta\lambda)$.
\end{theorem}
\begin{proof}
\leanok
Compute the two utilities in each comparison and solve the linear inequality in $p$.
\end{proof}

\section{Temptation}
\label{sec:beh-gp}

\begin{theorem}[Gul--Pesendorfer: set betweenness and commitment]
\label{thm:beh-gp}
\lean{Econ.Temptation.W, Econ.Temptation.setBetweenness, Econ.Temptation.commitment_strict, Econ.Temptation.W_pair_succumb, Econ.Temptation.W_le_max_u, Econ.Temptation.W_union_of_not_tempting}
\leanok
For $W(A)=\max_A(u+v)-\max_Av$: $W(A)\ge W(B)$ implies $W(A)\ge W(A\cup B)\ge W(B)$; a resisted temptation makes
$\{x\}$ strictly better than $\{x,y\}$.
\end{theorem}
\begin{proof}
\leanok
$W(A\cup B)=\max(a,b)-\max(\alpha,\beta)$; check the four cases of which menu attains each maximum.
\end{proof}

\section{Logit and rational inattention}
\label{sec:beh-ri}

\begin{theorem}[Generalized logit solves the Gibbs problem]
\label{thm:beh-gibbs}
\lean{Econ.Logit.generalizedLogit_optimal, Econ.Logit.eq_generalizedLogit_of_optimal, Econ.Logit.objective_generalizedLogit, Econ.Logit.softmax_optimal}
\leanok
For $\lambda>0$ and $P>0$, $p\mapsto\sum p_iv_i-\lambda\,\mathrm{KL}(p\|P)$ is uniquely maximized by
$p_i=P_ie^{v_i/\lambda}/\sum_jP_je^{v_j/\lambda}$, with value $\lambda\log\sum_jP_je^{v_j/\lambda}$.
\end{theorem}
\begin{proof}
\leanok
Term by term, $\log x\le x-1$ with $x=q_i/p_i$ gives $p_iv_i-\lambda p_i\log(p_i/P_i)\le\lambda(q_i-p_i)+\lambda p_i\log Z$,
strictly unless $p_i=q_i$; sum over $i$.
\end{proof}

\begin{theorem}[Mutual information is the minimal average divergence]
\label{thm:beh-mi}
\lean{Econ.RationalInattention.avgKL_eq, Econ.RationalInattention.mutualInfo_le_avgKL, Econ.RationalInattention.kl_nonneg, Econ.RationalInattention.kl_eq_zero_iff, Econ.RationalInattention.marginal_isProb}
\leanok
\uses{thm:beh-gibbs}
$\sum_\omega g_\omega\,\mathrm{KL}(p_\omega\|Q)=I(p)+\mathrm{KL}(P\|Q)\ge I(p)$, $P$ the marginal.
\end{theorem}
\begin{proof}
\leanok
State by state the difference is $\sum_ig_\omega p_{\omega i}(\log P_i-\log Q_i)$; swap sums. Gibbs' inequality
is the Gibbs problem with $v=0$.
\end{proof}

\begin{lemma}[Gibbs problem with a vanishing default]
\label{lem:beh-gibbs0}
\lean{Econ.RationalInattention.gl_nonneg, Econ.RationalInattention.gl_eq_zero, Econ.RationalInattention.gl_isProb, Econ.RationalInattention.log_gl, Econ.RationalInattention.term_le0, Econ.RationalInattention.term_lt0, Econ.RationalInattention.objective_le0, Econ.RationalInattention.objective_gl0, Econ.RationalInattention.eq_gl_of_optimal0, Econ.RationalInattention.kl_nonneg0, Econ.RationalInattention.Z_pos_of_isProb}
\leanok
\uses{thm:beh-gibbs}
For $\lambda>0$ and a probability default $P\ge0$ (zeros allowed), the generalized logit $q$ is a probability vector
vanishing where $P$ does, and every $p\ll P$ satisfies $\sum p_iv_i-\lambda\,\mathrm{KL}(p\|P)\le\lambda\log Z$, with
equality only at $p=q$. Gibbs' inequality $\mathrm{KL}(p\|q)\ge0$ holds whenever $p\ll q$.
\end{lemma}
\begin{proof}
\leanok
Terms with $P_i=0$ have $p_i=q_i=0$ and contribute $0$. On the others $\log q_i=\log P_i+v_i/\lambda-\log Z$ and
$\log x\le x-1$ (strict for $x\ne1$) bounds each term as in Theorem~\ref{thm:beh-gibbs}. $Z>0$ because some $P_i>0$.
\end{proof}

\begin{theorem}[Mat\v{e}jka--McKay, Theorem 1]
\label{thm:beh-mm}
\lean{Econ.RationalInattention.optimal_is_generalizedLogit, Econ.RationalInattention.optimal_is_softmax_of_uniform, Econ.RationalInattention.sum_objective_eq, Econ.RationalInattention.rule_eq_zero_of_marginal_eq_zero}
\leanok
\uses{thm:beh-mi, thm:beh-gibbs, lem:beh-gibbs0}
If $p$ maximizes $\sum_\omega g_\omega\sum_ip_{\omega i}v_{\omega i}-\lambda I(p)$ with marginal $P$
(actions with $P_i=0$ allowed), then $p_{\omega i}=P_ie^{v_{\omega i}/\lambda}/\sum_jP_je^{v_{\omega j}/\lambda}$ in every state; with a
uniform marginal this is the multinomial logit.
\end{theorem}
\begin{proof}
\leanok
Actions with $P_i=0$ have $p_{\omega i}=0$ in every state, so $p_\omega\ll P$. The Gibbs bound is redone
for defaults that may vanish (with $p\ll P$, terms with $P_i=0$ are $0$ on both sides). Replace
$p_{\omega_0}$ by the generalized logit for default $P$; its marginal still vanishes where $P$ does, so
$\mathrm{KL}(P'\|P)\ge0$. The objective of $p$ equals the per-state Gibbs
objectives with default $P$; that of the deviation is at least its per-state Gibbs objectives
(Theorem~\ref{thm:beh-mi}). By uniqueness in the Gibbs problem the deviation is strictly better unless
$p_{\omega_0}$ already is the generalized logit.
\end{proof}

\section{Ambiguity}
\label{sec:beh-ambiguity}

\begin{theorem}[Ellsberg]
\label{thm:beh-ellsberg}
\lean{Econ.Ambiguity.ellsberg_not_seu, Econ.Ambiguity.ellsberg_meu, Econ.Ambiguity.ellsbergPriors_nonempty}
\leanok
No weights $\pi$ on \{red, black, yellow\} rank the red bet above black and black-or-yellow above
red-or-yellow. Maxmin EU with priors $(1/3,0,2/3)$ and $(1/3,2/3,0)$ does.
\end{theorem}
\begin{proof}
\leanok
The two strict inequalities say $\pi_R>\pi_B$ and $\pi_B>\pi_R$. The maxmin values are $1/3,0,1/3,2/3$.
\end{proof}

\begin{theorem}[Gilboa--Schmeidler axioms hold for maxmin EU]
\label{thm:beh-gs}
\lean{Econ.Ambiguity.meu, Econ.Ambiguity.meu_mono, Econ.Ambiguity.meu_mix_ge, Econ.Ambiguity.meu_mix_const, Econ.Ambiguity.meu_le_eu, Econ.Ambiguity.exists_eq_meu}
\leanok
$V(f)=\min_{\pi\in C}E_\pi f$ is monotone, concave under mixing (uncertainty aversion), and
$V(\alpha f+(1-\alpha)c)=\alpha V(f)+(1-\alpha)c$ for constants $c$, $\alpha\ge0$.
\end{theorem}
\begin{proof}
\leanok
Each $E_\pi$ is linear; a minimum of linear functions is concave and commutes with adding constants.
\end{proof}

\begin{definition}[Maxmin functional and prior set]
\label{def:beh-maxmin-functional}
\lean{Econ.Ambiguity.IsMaxminFunctional, Econ.Ambiguity.priorSet, Econ.Ambiguity.priors, Econ.Ambiguity.IsMaxminFunctional.map_zero, Econ.Ambiguity.IsMaxminFunctional.map_const, Econ.Ambiguity.IsMaxminFunctional.smul, Econ.Ambiguity.IsMaxminFunctional.add_neg_le}
\leanok
A functional $I$ on utility acts $f:S\to\mathbb{R}$ is a maxmin functional if it is monotone, $I(f+c)=I(f)+c$,
$I(tf)=tI(f)$ for $t>0$, and $I(f)+I(g)\le I(f+g)$. Its prior set is
$C_I=\{p\in\Delta(S): E_pg\ge I(g)\text{ for all }g\}$.
\end{definition}

\begin{lemma}[Supporting expectation]
\label{lem:beh-gs-support}
\lean{Econ.Ambiguity.IsMaxminFunctional.exists_support}
\leanok
\uses{def:beh-maxmin-functional}
For a maxmin functional $I$ and any act $f_0$ there is $p\in C_I$ with $E_pf_0=I(f_0)$.
\end{lemma}
\begin{proof}
\leanok
$N(g)=-I(-g)$ is sublinear. On the line through $f_0\ne0$ the functional $tf_0\mapsto tI(f_0)$ lies below $N$:
for $t>0$ because $I(f_0)+I(-f_0)\le I(0)=0$, for $t<0$ with equality. Hahn--Banach extends it to a linear $L\le N$,
i.e.\ $L\ge I$. Then $p_s=L(e_s)\ge I(e_s)\ge I(0)=0$ and $L(1)=1$ from $I(\pm1)=\pm1$. For $f_0=0$ use any such $p$
from the line through the constant $1$.
\end{proof}

\begin{theorem}[Gilboa--Schmeidler Lemma 3.5: functional form]
\label{thm:beh-gs-functional}
\lean{Econ.Ambiguity.maxmin_of_isMaxminFunctional, Econ.Ambiguity.isMaxminFunctional_of_isLeast, Econ.Ambiguity.eu_add_smul, Econ.Ambiguity.continuous_eu, Econ.Ambiguity.convex_priors, Econ.Ambiguity.isClosed_priors, Econ.Ambiguity.eu_const_add, Econ.Ambiguity.eu_smul', Econ.Ambiguity.eu_add', Econ.Ambiguity.eu_mono', Econ.Ambiguity.isMaxminFunctional_meu}
\leanok
\uses{lem:beh-gs-support}
$I$ is a maxmin functional iff $I(f)=\min_{p\in C}E_pf$ for a set of priors $C$ at which the minimum is attained;
$C_I$ is then nonempty, closed and convex.
\end{theorem}
\begin{proof}
\leanok
Lemma~\ref{lem:beh-gs-support} gives attainment; $C_I$ is an intersection of closed half-spaces with the simplex.
Conversely each property passes from the linear $E_p$ to the minimum, using the minimizer at the relevant act.
\end{proof}

\begin{definition}[Gilboa--Schmeidler axioms]
\label{def:beh-gs-axioms}
\lean{Econ.Ambiguity.GSAxioms, Econ.Ambiguity.constAct, Econ.Ambiguity.mixAct, Econ.Ambiguity.StrictPref, Econ.Ambiguity.mixAct_const}
\leanok
On utility acts: weak order; C-independence ($f\succsim g$ iff $\alpha f+(1-\alpha)c\succsim\alpha g+(1-\alpha)c$ for
constants $c$, $\alpha\in(0,1)$); Archimedean; monotonicity; uncertainty aversion ($f\sim g$ implies
$\alpha f+(1-\alpha)g\succsim f$); nondegeneracy normalized as $c_1\succ c_0$.
\end{definition}

\begin{theorem}[Gilboa--Schmeidler 1989, Theorem 1]
\label{thm:beh-gs-main}
\lean{Econ.Ambiguity.gilboa_schmeidler, Econ.Ambiguity.gilboa_schmeidler_unique, Econ.Ambiguity.eq_priorSet, Econ.Ambiguity.gsAxioms_of_isMaxminFunctional, Econ.Ambiguity.gsAxioms_ellsberg, Econ.Ambiguity.GSAxioms.refl, Econ.Ambiguity.GSAxioms.nonempty, Econ.Ambiguity.GSAxioms.const_le_iff, Econ.Ambiguity.GSAxioms.strict_of_not, Econ.Ambiguity.GSAxioms.exists_ce, Econ.Ambiguity.GSAxioms.ce, Econ.Ambiguity.GSAxioms.ce_spec, Econ.Ambiguity.GSAxioms.ce_eq, Econ.Ambiguity.GSAxioms.rep, Econ.Ambiguity.GSAxioms.ce_const, Econ.Ambiguity.GSAxioms.ce_mix_const, Econ.Ambiguity.GSAxioms.ce_smul_lt_one, Econ.Ambiguity.GSAxioms.ce_pos_hom, Econ.Ambiguity.GSAxioms.ce_const_add, Econ.Ambiguity.GSAxioms.isMaxminFunctional_ce}
\leanok
\uses{def:beh-gs-axioms, thm:beh-gs-functional}
A preference satisfies the axioms iff it is represented by maxmin expected utility over a nonempty closed convex
set of priors, and that set is unique (necessity holds for any maxmin functional). Acts are utility acts with
values in all of $\mathbb{R}^S$. Maxmin EU over the two Ellsberg priors satisfies the axioms.
\end{theorem}
\begin{proof}
\leanok
Constants are ranked by value (C-independence moves any strict pair to $(0,\delta)$, then to $(0,1)$). The Archimedean
axiom gives a certainty equivalent $I(f)$: the supremum of constants below $f$ is indifferent to $f$. C-independence
with the constant $0$ gives homogeneity on $(0,1)$, hence for all $t>0$; with $\alpha=1/2$ it gives constant
additivity. Uncertainty aversion on $f-I(f)\sim g-I(g)$ gives superadditivity. Apply
Theorem~\ref{thm:beh-gs-functional}. Uniqueness: two representations have the same certainty equivalents, and a closed convex set representing $I$ is
$C_I$, since a prior outside it is strictly separated from it by some act $g$, giving $E_pg<\min_CE_qg=I(g)$. For necessity, the Archimedean step uses concavity
($\alpha$ near $1$) and the 1-Lipschitz bound $I(h+\beta(f-h))\le I(h)+\beta\sum_s|f_s-h_s|$ ($\beta$ near $0$).
\end{proof}

\section{Stochastic choice}
\label{sec:beh-stochastic-choice}

\begin{definition}[Choice rule, Luce rule]
\label{def:beh-luce}
\lean{Econ.StochasticChoice.IsChoiceRule, Econ.StochasticChoice.luce, Econ.StochasticChoice.probSet, Econ.StochasticChoice.ChoiceAxiom}
\leanok
A stochastic choice rule assigns each nonempty menu $A$ a probability on $A$. The Luce rule with positive weights
is $P(a,A)=w_a/\sum_{b\in A}w_b$. The choice axiom: $P(a,A)=P(B,A)P(a,B)$ for $a\in B\subseteq A$.
\end{definition}

\begin{theorem}[Luce 1959]
\label{thm:beh-luce}
\lean{Econ.StochasticChoice.luce_isChoiceRule, Econ.StochasticChoice.luce_iia, Econ.StochasticChoice.luce_choiceAxiom, Econ.StochasticChoice.eq_luce_of_choiceAxiom, Econ.StochasticChoice.softmax_eq_luce}
\leanok
\uses{def:beh-luce}
Luce rules are choice rules satisfying IIA and the choice axiom. Conversely a choice rule with full support on every
menu that satisfies the choice axiom is the Luce rule with $w_a=P(a,X)$. Logit is Luce with $w=e^{v/\lambda}$.
\end{theorem}
\begin{proof}
\leanok
Direct computation. For the converse apply the axiom with $B=A\subseteq X$: $P(a,X)=P(A,X)P(a,A)$ and
$P(A,X)=\sum_{b\in A}w_b$.
\end{proof}

\begin{definition}[Random utility]
\label{def:beh-rum}
\lean{Econ.StochasticChoice.Ranking, Econ.StochasticChoice.IsTop, Econ.StochasticChoice.IsRUM, Econ.StochasticChoice.rumP, Econ.StochasticChoice.bm, Econ.StochasticChoice.isRUM_pointMass}
\leanok
A RUM is a probability $\mu$ on strict rankings; $P(a,A)$ is the probability that $a$ is best in $A$.
The Block--Marschak polynomial is $\mathrm{BM}(a,A)=\sum_{B\supseteq A}(-1)^{|B\setminus A|}P(a,B)$.
\end{definition}

\begin{theorem}[Random utility: regularity and Block--Marschak]
\label{thm:beh-bm}
\lean{Econ.StochasticChoice.sum_isTop, Econ.StochasticChoice.rumP_isChoiceRule, Econ.StochasticChoice.rumP_regularity, Econ.StochasticChoice.isTop_union, Econ.StochasticChoice.bm_rank_term, Econ.StochasticChoice.bm_rumP, Econ.StochasticChoice.bm_rumP_nonneg}
\leanok
\uses{def:beh-rum}
A RUM is a choice rule satisfying regularity, and $\mathrm{BM}(a,A)$ equals the probability that $a$ is best in $A$
and everything outside $A$ is ranked above $a$; in particular $\mathrm{BM}\ge0$ (Block--Marschak 1960).
\end{theorem}
\begin{proof}
\leanok
Fix a ranking. For $B=A\cup C$ with $C\subseteq X\setminus A$, $a$ is best in $B$ iff it is best in $A$ and $C$
avoids the alternatives outside $A$ ranked above $a$. So the alternating sum runs over subsets of
$D=\{c\notin A: a\text{ above }c\}$ and equals $[D=\emptyset]$. Average over rankings.
\end{proof}
